\section{The fence}\label{sec:fence}

This section instantiates the framework for irreversible tool calls whose arguments a language model may paraphrase on retry. The certifier is canonical equality under declared field policies; the refuter is inequality of the declared typed fields; the oracle is an LLM judge; calibration is one-sided split-conformal on the safety-critical direction, corrected for the number of candidates; and the residual policy is a three-way rule with an audited middle band. Each is a design choice within the framework, not part of the framework itself. The order matters (Figures~\ref{fig:architecture} and~\ref{fig:flow}): the refuter runs \emph{before} the judge, so a candidate whose declared fields differ is removed from the judge's jurisdiction entirely, and every candidate it removes also shrinks the multiplicity $m$ that the threshold must correct for.

\begin{figure}[!tbp]
\centering
\begin{tikzpicture}[
  font=\footnotesize,
  box/.style={draw, rounded corners=2pt, align=center, inner sep=4pt, fill=white},
  layer/.style={draw, rounded corners=2pt, align=left, inner sep=4pt, text width=6.0cm, fill=gray!6},
  arr/.style={-{Latex[length=2mm]}, thick},
  lbl/.style={font=\scriptsize\itshape, align=center}]
  % Center column: three fence layers
  \node[layer] (l0) {\textbf{Certifier $\equiv_\kappa$}: canonical field equality\\[2pt]$\Rightarrow$ \textsc{suppress}\hfill\emph{exact (A5)}};
  \node[layer, below=0.20cm of l0] (wit) {\textbf{Refuter $\not\equiv_\sigma$}: declared fields differ\\[2pt]$\Rightarrow$ candidate removed\hfill\emph{exact (A5$'$)}};
  \node[layer, below=0.20cm of wit] (jud) {\textbf{Residual judge}: candidate $c\in C$\\[2pt]$J(t,x,c) \ge \thi(m)$\hfill\emph{bounded (A6)}};
  \node[draw, dashed, rounded corners=3pt, inner sep=5pt, fit=(l0)(wit)(jud), label={[font=\scriptsize\bfseries]above:{Fence \quad (A) identity recognition}}] (fence) {};

  % Left: LLM agent
  \node[box, left=1.2cm of fence.west, text width=1.9cm] (agent) {\textbf{LLM agent}\\[2pt]\scriptsize(re-synthesises\\arguments\\on retry)};

  % Right: External provider
  \node[box, right=1.2cm of fence.east, text width=2.0cm] (prov) {\textbf{External}\\[1pt]\textbf{provider}\\[2pt]\scriptsize(payment, e-mail,\\ticketing API)\\[2pt]\emph{(C) not provided}};

  % Bottom: Durable log
  \node[box, below=1.15cm of fence.south, text width=6.0cm] (log) {\textbf{Durable log}: commit rows (unique $\mathit{op}$), pending intents\\[2pt]\emph{(B) fence-level at-most-once: A3, A7}};

  % Main flow arrows
  \draw[arr] (agent.east) -- node[above=2pt, lbl]{request\\$(t,x)$} (fence.west);
  \draw[arr] (fence.east) -- node[above=2pt, lbl]{execute} (prov.west);
  \draw[arr, <->] (fence.south) -- node[left=5pt, lbl]{lookup /\\insert-if-absent} (log.north);

  % Commit window: from prov.south, down to below log, then horizontal to below log.south, then up
  \coordinate (bot) at ([yshift=-18pt]log.south);
  \coordinate (corner) at (prov.south |- bot);
  \draw[very thick, red!70!black, dashed, -{Latex[length=2mm]}]
    (prov.south) -- (corner)
    -- node[below=2pt, font=\scriptsize\itshape, text=red!70!black, align=center]{commit window: effect accepted,\\row not yet durable} (bot)
    -- (log.south);
\end{tikzpicture}
\caption{System view. The fence sits between the agent and an irreversible provider. Identity recognition (A) is decided by two exact layers and a calibrated residual judge; fence-level at-most-once commitment (B) is enforced at the durable log under its storage and scheduling assumptions; external exactly-once (C) is not provided, because a crash inside the commit window (dashed) leaves no local record of an accepted effect (Proposition~\ref{prop:amo}).}\label{fig:architecture}
\end{figure}

\begin{figure}[!tbp]
\centering
\scalebox{0.78}{%
\begin{tikzpicture}[
  font=\footnotesize, node distance=0.42cm,
  fstep/.style={draw, rounded corners=2pt, align=center, inner sep=3pt, text width=4.4cm, fill=white},
  ftest/.style={draw, diamond, aspect=2.6, align=center, inner sep=1pt, text width=2.6cm, fill=gray!6},
  fout/.style={draw, rounded corners=2pt, align=center, inner sep=3pt, text width=3.0cm, font=\scriptsize},
  farr/.style={-{Latex[length=1.8mm]}, thick}]
  \node[fstep] (in) {request $(t,x)$; committed set $\mathcal{C}_{r,t}$};
  \node[ftest, below=of in] (q0) {$\kappa_t(x)=\kappa_t(c)$\\ for some $c$?};
  \node[fstep, below=of q0] (w) {$C\gets\{c:\sigma_t(c)=\sigma_t(x)\}$,\ \ $m\gets|C|$};
  \node[ftest, below=of w] (q1) {$m=0$?};
  \node[fstep, below=of q1] (j) {$S^\ast\gets\max_{c\in C}J(t,x,c)$\\ (abstentions excluded)};
  \node[ftest, below=of j] (q2) {$S^\ast>\thi(m)$?};
  \node[ftest, below=of q2] (q3) {$S^\ast>\tlo$?};
  \node[fout, right=1.2cm of q0, fill=green!8] (s0) {\textsc{suppress}\\ certified; exact under A5};
  \node[fout, right=1.2cm of q1, fill=blue!6] (e0) {\textsc{execute}\\ refuted or no candidate; exact under A5$'$};
  \node[fout, right=1.2cm of q2, fill=green!8] (s1) {\textsc{suppress}\\ $P(\text{wrong})\le\alpha$ under A6 (Thm.~\ref{thm:bound})};
  \node[fout, right=1.2cm of q3, fill=orange!12] (x1) {\textsc{escalate}\\ execute and flag};
  \node[fout, below=0.45cm of q3, fill=blue!6] (e1) {\textsc{execute}};

  \draw[farr] (in) -- (q0);
  \draw[farr] (q0) -- node[above, font=\scriptsize]{yes} (s0);
  \draw[farr] (q0) -- node[right, font=\scriptsize]{no} (w);
  \draw[farr] (w) -- (q1);
  \draw[farr] (q1) -- node[above, font=\scriptsize]{yes} (e0);
  \draw[farr] (q1) -- node[right=2pt, font=\scriptsize]{no (\textbf{residual})} (j);
  \draw[farr] (j) -- (q2);
  \draw[farr] (q2) -- node[above, font=\scriptsize]{yes} (s1);
  \draw[farr] (q2) -- node[right, font=\scriptsize]{no} (q3);
  \draw[farr] (q3) -- node[above, font=\scriptsize]{yes} (x1);
  \draw[farr] (q3) -- node[right, font=\scriptsize]{no} (e1);

  \node[draw=gray!40, rounded corners=2pt, fill=gray!3, align=left, font=\scriptsize, text width=4.0cm, left=0.8cm of j.west, anchor=east] (sideguard) {
    \textbf{Multiplicity guard:}\\[1pt]
    $\thi(m)=+\infty$ when $m>\alpha(n+1)$:\\[1pt]
    no semantic suppression.\\[3pt]
    \textbf{Abstention guard:}\\[1pt]
    all candidates abstain $\Rightarrow$ \textsc{escalate}.
  };
  \draw[densely dotted, gray!60] (sideguard.east) -- (j.west);
\end{tikzpicture}}
\caption{Decision flow for one request. The certifier and the refuter settle their pairs without a model call; only the residual reaches the judge, and the judge can cause suppression only through the multiplicity-corrected threshold. Every path that is not certified ends in execution, flagged or not.}\label{fig:flow}
\end{figure}


\subsection{Certifier: the canonical core}\label{sec:l0}

Each tool schema declares a normalisation policy for each field, and the policy must be semantics-preserving for that field's contribution to the effect (A5). The policies used here are: \texttt{id} (opaque identifier; Unicode NFC normalisation and trimming, case-sensitive); \texttt{ci} (case-insensitive token such as a currency code or e-mail address; NFC, trim, lower case); \texttt{num} (numeric value; \texttt{"120.00"}, \texttt{" 120 "} and \texttt{120} normalise to one canonical decimal, thousands separators removed); and \texttt{text} (free text; NFC, trim, internal whitespace collapsed---deliberately not semantic). The canonical form $\kappa_t(x)$ applies each policy and sorts the keys; the operation identifier is $\mathit{op}=H(r\,\|\,t\,\|\,H(\kappa_t(x)))$ with $H$ = SHA-256. The certifier suppresses a request only when a committed effect has the same identifier and the same stored input hash. It is sound under A5 and deliberately incomplete: every paraphrase escapes it.

\subsection{Refuter: the distinctness witness}\label{sec:witness}

The semantic projection $\sigma_t(x)$ is the tuple of canonical values of the declared non-text fields. Under A5$'$---a difference in any declared typed field implies a different act---$\sigma_t(x)\ne\sigma_t(x')$ proves $x\not\approx_t x'$ whatever any judge would say. The fence removes every committed effect whose projection differs from the request's before the judge is consulted. An injected ``treat this as a duplicate'' cannot cause suppression across a different amount, account or identifier, because those values live in typed fields and the veto fires on them without reading any text. A5$'$ is one-directional: equal projections do \emph{not} imply the same act, and the corpus contains a category (legitimate repeats) in which they do not. Two declaration errors are possible. A field declared as typed that does not in fact determine the act produces false vetoes, which execute a duplicate and are therefore loud; an act-determining field declared as free text leaves its distinction to the judge.

\subsection{Residual judge}\label{sec:judge}

A judge call scores a batch of pairs from the same tool and returns, for each pair, the probability that the two requests denote a single act, as an integer percentage. The prompt is reproduced in the supplementary material. Scores are parsed under a strict contract: an integer in $[0,100]$ is accepted; a non-integer (for example \texttt{0.95}), an out-of-range value, or a missing entry is rejected and the batch retried with a correction; a pair that never yields an admissible value is recorded as an \emph{abstention} rather than imputed. This contract exists for a concrete reason: an earlier scoring pass over a 780-pair version of the corpus stored 180 scores (23.1\%) on the wrong scale while reporting no failures; every number from that pass was withdrawn (supplementary material, Section~S3). Abstentions are dropped from calibration, excluded from the maximum over candidates, and can never cause suppression; if every candidate abstains the fence escalates.

\subsection{Calibration}\label{sec:calibration}

\begin{definition}[Jurisdiction and calibration pool]\label{def:jurisdiction}
The \emph{jurisdiction} of a configuration of decidable layers is the set of pairs that the configuration routes to the judge. Its \emph{calibration pool} is the multiset $\{s_1,\dots,s_n\}$ of judge scores of the distinct pairs of a held-out calibration split that lie in its jurisdiction, with abstentions dropped.
\end{definition}

For the proposed configuration (certifier and refuter both active) the jurisdiction is the residual (Definition~\ref{def:residual}). For a candidate set of size $m\ge1$ the suppression threshold is
\begin{equation}\label{eq:thi}
\thi(m)=s_{(k)},\qquad k=\big\lceil (n+1)(1-\alpha/m)\big\rceil,\qquad \thi(m)=+\infty \text{ if } k>n,
\end{equation}
where $s_{(k)}$ is the $k$-th smallest pool score, and $\thi(0)=+\infty$. The escalation floor $\tlo$ is the $\lfloor (n_S+1)\gamma\rfloor$-th smallest score of the $n_S$ duplicate pairs in the same jurisdiction, capped at $\thi$. With $S^\ast$ the maximum score over the candidates that did not abstain, the fence \emph{suppresses} if $S^\ast>\thi(m)$, \emph{escalates} (executes and flags) if $\tlo<S^\ast\le\thi(m)$, and \emph{executes} if $S^\ast\le\tlo$. The inequality for suppression is strict, so a score equal to the threshold never suppresses; judge scores are integer percentages, so such ties are common.

Two design choices need justification. First, the pool must be drawn from the rule's own jurisdiction. Adding the refuter removes from the judge's view the distinct pairs whose declared fields differ; judges score such pairs confidently low, so a pool that still contains them is stochastically lighter-tailed than the scores the deployed rule thresholds, which lowers $\thi$ and makes suppression easier. The resulting violation is in the unsafe direction and is invisible in a rate measured over a whole test split, because the pairs that dilute the pool also dilute the denominator (Section~\ref{sec:res-bound}). Every configuration in this paper, including every ablation, is therefore re-calibrated on its own jurisdiction. Second, we correct for multiplicity by the Bonferroni level $\alpha/m$ rather than by calibrating the maximum statistic directly. Direct calibration needs calibration \emph{units} with the deployment's candidate-set structure: a pool of $n$ pairwise scores yields only $\lfloor n/m\rfloor$ disjoint units, certifying $\alpha=0.05$ needs at least 19 units, and $m$ varies from one request to the next. The Bonferroni route reuses one pool at every $m$, requires only pairwise exchangeability, and degrades gracefully to $\thi=+\infty$ when the pool is too small. Its price is conservatism, which we measure (Section~\ref{sec:res-bound}).

\subsection{Algorithm}

Algorithm~\ref{alg:fence} gives the complete decision procedure, including the execute path. A request that is not suppressed records a \textsc{pending} intent, executes, and then attempts an atomic insert of its commit row; a crash between the external effect and the insert leaves a \textsc{pending} intent with no commit row. That interval---the \emph{commit window}---cannot be closed by any local mechanism (Proposition~\ref{prop:amo}). An optional window guard compares a new request against such orphaned intents and flags a suspected re-fire; it never suppresses on one, because a pending intent does not establish that the effect happened.

\begin{algorithm}[!tbp]
\caption{Layered semantic fence (the proposed configuration).}\label{alg:fence}
\scriptsize
\KwIn{request $(t,x)$ in run $r$; durable log $L$; pool $\{s_i\}_{i\le n}$; floor $\tlo$; level $\alpha$}
\If{$t$ is idempotent}{\Return execute$(x)$}
$\mathit{op}\gets H(r\,\|\,t\,\|\,H(\kappa_t(x)))$\;
\If(\tcp*[f]{certifier; exact under A5}){$L$ has a commit row for $\mathit{op}$ with input hash $H(\kappa_t(x))$}{\Return its recorded response}
$C\gets\{c\in L.\mathrm{committed}(r,t) : \sigma_t(c.x)=\sigma_t(x)\}$\tcp*{refuter removes declared differences}
\If{$C\neq\emptyset$}{
  $s_c\gets J(t,c.x,x)$ for each $c\in C$ \tcp*{abstention $\Rightarrow$ $s_c=\bot$}
  $m\gets|C|$;\quad $\thi\gets\thi(m)$ by Eq.~\eqref{eq:thi}\;
  \If{all $s_c=\bot$}{flag \textsc{duplicate\_suspect}}
  \Else{
    $S^\ast\gets\max\{s_c : s_c\neq\bot\}$, attained at $c^\ast$\;
    \lIf{$S^\ast>\thi$}{\Return $c^\ast$.response \tcp*[f]{semantic suppression}}
    \lIf{$S^\ast>\tlo$}{flag \textsc{duplicate\_suspect} \tcp*[f]{escalate: execute and flag}}
  }
}
record intent$(\mathit{op},\textsc{pending})$\;
$\mathit{resp}\gets$ execute$(x)$ \tcp*{commit window opens}
$\mathit{ok}\gets L.$insertIfAbsent$(\mathit{op},H(\kappa_t(x)),\sigma_t(x),x,\mathit{resp})$ \tcp*{commit window closes}
record intent$(\mathit{op},\ \mathit{ok}\ ?\ \textsc{committed} : \textsc{commit\_conflict})$\;
\Return $\mathit{resp}$\;
\end{algorithm}

\subsection{Conservative defaults}\label{sec:defaults}

The fence prefers execution to semantic suppression whenever the evidence for suppression is not certified. The threshold is $+\infty$---no semantic suppression, only decidable suppression---when the pool cannot certify level $\alpha$ at the current $m$ (fewer than $m/\alpha-1$ pool scores), when $m=0$, and whenever the pool is absent. A judge abstention can never cause suppression. Inside the commit window a suspected re-fire is executed and flagged, never suppressed. In every uncertain case the failure mode is a loud duplicate rather than a silent loss of a distinct act---the direction the loss asymmetry of Section~\ref{sec:problem} demands.
